\documentclass{subfiles}
\begin{document}
    In this section, we first introduce the notion of group, ring and field,
        we then define what a field extension is and how to it is used. 
        In describe what above, we assume the reader has a basic understanding of 
        modulo arithmetics.

    \begin{definition*}[Group]
    Let \(G\) be a set, and let \(*\) denote a binary operation on \(G\)'s elements.
        We call \emph{group} the pair \((G, *)\) if the following hold:
        \begin{enumerate}
            \item There exists a neutral element in \(G\);
                that is, 
                \[
                    \exists e \in G \text{ s.t. } \forall a \in G, a * e = e * a = a
                \]
            \item Each element has an inverse, i.e, 
                \[
                    \forall a \in G, \exists a^{-1} \text{ s.t. } a * a^{-1} = a^{-1} * a = e 
                \]
            \item The operation \(*\) is associative, meaning that 
                \[
                    \forall a, b, c \in G, (a * b) * c = a * (b * c)
                \]
        \end{enumerate}
        Additionally, if \(*\) is such that \(a * b = b * a\), i.e., 
            \(*\) is commutative, we say that the group is \emph{abelian}.
    \end{definition*}
    The number of elements of a group \(G\) is said to be its order \(\abs{G}\).

    \begin{definition*}[Cyclic group]
        Let \((G, *)\) be a group, and let \(a \in G\). We denote by \(a^{n}\)
            the \(n\)-th power of \(a\), and denote the set of all such powers by 
            \(\gen{a}\). We say that \((G, *)\) is cyclic if \(\gen{a} = G\).
    \end{definition*}

    \begin{definition*}[Sub-group]
        Given \((G, *)\) a group, we say that \(H \subseteq G\) forms a sub-group
            under \(*\), if \((H, *)\) is itself a group.
    \end{definition*}

    The following is one of the core theorems in group theory,
        and its extentions to rings, comes in handy when dealing with fields.
        \begin{theorem*}[Lagrange's theorem]
            Let \((G, *)\) be a group and let \((h, *)\) be a subgroup of \(G\).
            Then, it holds that \(\abs{H} \text{ divides } \abs{G}\).
        \end{theorem*}

    Groups are one of several algebraic structure, among which we also consider 
        rings and fields. The latter will be at the center of our discussion on 
        cryptography.

    \begin{definition*}[Ring]
        Let \(R\) be a set, and let \(+, *\) be two distinct binary operations
            on \(R\)'s elements. We say that \((R, +, *)\) is a ring if 
            \begin{enumerate}
                \item \((R, +)\) is an abelian group;
                \item The operation \(*\) is associative;
                \item The operation \(*\) distibutes over \(+\); that is,
                    \[
                        \forall a, b, c \in G, a * (b + c) = a * b + a * c.
                    \]
            \end{enumerate}
        Additionally, if \(R\) has a multiplicative unit, i.e., \(\exists e' \in R\)
        such that \(e' * a = a * e' = a, \forall a \in R\) and \(*\) is commutative,
        we say that \((R, +, *)\) is a commutative ring with unity.
    \end{definition*}
    In what follows we denote by \(\field^{*}\) the ring \(\field \setminus \set{0}\).

    \begin{definition*}[Field]
        Given a commutative ring with unity \((R, +, *)\), we say that \((R, +, *)\)
            is a field if \((R^{*}, +, *)\) is commutative.
    \end{definition*}
    Similarly for groups, each field has an order given by the number of its elements.
        In these notes we mostly consider fields of order a prime and denote these by \(\gf{p}\),
        for some prime number \(p\).
\end{document}
